\documentclass[11pt,a4paper]{article}
\usepackage[margin=2.5cm]{geometry}
\usepackage{amsmath,amssymb,booktabs,graphicx,hyperref,natbib}
\usepackage[T1]{fontenc}
\usepackage{float}
\hypersetup{colorlinks=true,linkcolor=blue,citecolor=blue,urlcolor=blue}
\title{\textbf{Energy Shocks and the Green Adoption Channel in a Small Open Economy}}
\author{Guillaume, Boris Chafwehe, Francesca Diluiso}
\date{\today}

\begin{document}
\maketitle

\begin{abstract}
\noindent We add an extensive-margin brown/green durable adoption choice to the
small open economy HANK model of \citet{ARS2024}. Households choose discretely
whether to switch to a green durable that insulates them from fossil energy
prices, paying a resource cost. We study brown energy \emph{price} and
\emph{supply} shocks under the two fiscal responses studied in the recent
literature: a retail price cap \citep{Langot2026} and a Slutsky-compensating
transfer indexed to pre-crisis energy consumption \citep{Bayer2026}. The price
cap stabilises output but mechanically eliminates the green adoption margin,
while the transfer delivers comparable stabilisation and leaves adoption
intact. Shielding consumers from the relative price is therefore not neutral
for the energy transition. This is our contribution and, we think, its natural
boundary: the three closest papers---\citet{ARS2024}, \citet{Bayer2026} and
\citet{Langot2026}---all hold the household's energy technology fixed, and we
add the one margin they omit rather than the two-country or labour-market
machinery they contain. The transition loss is linear in the cap's intensity
under a price shock and convex under a supply shock. As a secondary
observation, the \emph{sign} of the cap's private-welfare effect flips with the
elasticity of energy supply, matching the sign each of \citet{Langot2026} and
\citet{Bayer2026} obtains in its own closure; we read this as locating one axis
of their contrast, not as nesting their models. The policy \emph{ordering} is
robust to calibration; the \emph{magnitude} of the adoption channel, and even
the sign of its output effect under a price shock, depend on the calibrated
switching cost and taste scale, which we report and discuss.
\end{abstract}

\section{Introduction}

Three recent papers study fiscal responses to the 2022 European energy crisis
in HANK frameworks. \citet{ARS2024} study a small open economy and stress the
negative externality when all importers subsidise simultaneously.
\citet{Bayer2026} use a two-country HANK with \emph{fixed} euro-area energy
supply and find subsidies are beggar-thy-neighbour while Slutsky transfers are
not. \citet{Langot2026} evaluate the French \emph{bouclier tarifaire} assuming
\emph{perfectly elastic} energy supply and find the price cap performs well.
All three treat the household's energy technology as fixed.

We relax exactly that margin, and only that margin. Households can adopt a
green durable that insulates them from the fossil price, at a resource cost.
The energy shock raises the return to adoption, so the crisis itself can
accelerate the transition---unless policy removes the price signal that drives
it. This is the contribution: an endogenous green-adoption margin, absent from
all three references. We deliberately do \emph{not} try to nest the two-country
spillovers of \citet{Bayer2026} or the labour-cost employment channel of
\citet{Langot2026}. A small open economy with an exogenous foreign block
contains neither, and competing on that machinery is not the point; our single
comparative advantage is the margin the literature leaves out, and we play only
that.

Because the price-versus-supply distinction turns out to matter for the
adoption margin, we do let the energy-supply elasticity vary: with
$\varepsilon^{E}=\infty$ the economy is a price taker (elastic supply, as in
Langot); with $\varepsilon^{E}$ finite the quantity is fixed and the world
price clears the market (as in Bayer). Both closures are reported below. This
is a modelling convenience for our own question, not a claim to reconcile the
two papers: the supply elasticity is one dimension along which they differ, and
not the one that carries their headline mechanisms (Section~\ref{sec:lit}).

\section{Model}

\subsection{Environment}

The aggregate structure is \citet{ARS2024}: a small open economy with sticky
wages, a CES consumption basket nesting energy against a home/foreign bundle,
imported energy and foreign goods with sticky retail prices, and a mutual fund
holding claims on domestic firms, importers and a share $\zeta^{E}$ of the
energy endowment. We build the model literally on the ARS aggregate blocks and
replace its representative, homothetic household with a heterogeneous-agent
household that carries a discrete durable state; everything else in the
aggregate economy is inherited unchanged.

\subsection{Model at a glance: the maintained assumptions}
\label{sec:assumptions}

Because the household block departs from ARS in several coupled ways, we list
the maintained assumptions explicitly before deriving anything. A reader who
accepts this list can reconstruct every result.

\begin{enumerate}\itemsep2pt
\item \textbf{One new margin.} The only structural addition to \citet{ARS2024}
is a discrete brown/green durable adoption choice. Conditional on the durable
state, preferences are ARS's: the same CES energy nest ($\alpha_E$, $\eta_E$)
and the same home/foreign nest. There is no within-durable-type
energy-share heterogeneity.
\item \textbf{Green insulates from the fossil price.} A green adopter buys
energy at $P^{E}_{G}$, which in the baseline ($\rho_g=0$) is constant and does
\emph{not} respond to the fossil shock. Adoption is therefore, in the baseline,
full insulation---an upper bound on the channel.
\item \textbf{Green is resource-cheaper; brown is imported.} The energy an
adopter uses comes from a technology whose marginal cost is near zero once the
durable is installed---photovoltaic electricity, or the ambient heat a heat
pump moves per unit of electricity (its coefficient of performance)---whereas
brown energy is imported. How this physical saving enters the balance of
payments is a modelling \emph{choice}, made explicit and varied in
Section~\ref{sec:bop}, not an identity: under the baseline \emph{import} booking
all energy is imported at the world brown price and only the adopter's price
saving $(P^{E}_{B}-P^{E}_{G})\,C^{G}_{E}$ is netted off as an external import
reduction, so green remains partly in the import bill and the switching cost
$\psi_g D^{sw}$ is itself booked as an import; under the \emph{domestic} booking
green energy and installation are domestic value added and only brown is
imported. Whether the switch reduces net imports \emph{on balance} is therefore
not assumed here---it is the object of Section~\ref{sec:signmap}. We do not
model the technology that supplies the green energy (no green production sector,
no green investment).
\item \textbf{The switching cost is a resource cost; under the baseline it is
booked externally.} Adopting costs $\psi_g$ units, entered in the current
account as an import under the baseline import booking---by accounting necessity
given ARS's markup/capitalised-profits convention, since booking it as domestic
demand would create a dividend owned by nobody. The domestic booking instead
rebates it as domestic value added (Section~\ref{sec:bop}).
\item \textbf{Unit of account is the domestic good.} We set $p^{num}=p_H$
(Section~\ref{sec:num}). This is a measurement choice, orthogonal to the
nominal anchor; it keeps the resource constraint independent of how the CPI
weights brown and green energy.
\item \textbf{Energy supply is a single dial.} $\varepsilon^{E}=\infty$ gives
the price-taking closure (exogenous world price, price shock);
$\varepsilon^{E}$ finite fixes the quantity and lets the world price clear
(supply shock). The steady state is identical under both.
\item \textbf{Monetary baseline is the constant-real-rate rule} of
\citet{ARS2024}. Under this rule the real exchange rate is pinned, so a
domestic transfer has no terms-of-trade effect; we also report a Taylor rule.
\item \textbf{First order throughout.} All impulse responses are first-order
(sequence-space Jacobian). The steady state is common to every experiment,
since both policy instruments are zero at the steady state.
\end{enumerate}

\subsection{Unit of account}
\label{sec:num}

The household block is numeraire-generic: it reads a price of the unit of
account relative to the CPI, $p^{num}$, a real return
$1+r^{num}=(1+r)\,p^{num}_{-1}/p^{num}$, and after-tax labour income
$atw^{num}=atw/p^{num}$, and returns asset holdings converted by
$A^{CPI}=A\,p^{num}$ before meeting CPI-denominated objects in asset market
clearing and the current account. We set $p^{num}=p_{H}$, the domestic good.
This makes the resource constraint independent of how the CPI weights brown
and green energy, which matters once the two coexist (Section~\ref{sec:cpi}).
Everything contractually written in CPI units---fiscal transfers, the ARS
subsistence term, and the switching cost $\psi_g$---is divided by $p^{num}$
on the way into the budget.

Two remarks make the convention transparent. First, the choice of numeraire is
the choice of nominal anchor, not of deflator: the Fisher equation is written
once in CPI terms ($1+r^{ante}=(1+i)/(1+\pi)$), and $p^{num}=p_H$ then delivers,
after the relative-price conversion, a household discount factor deflated by
\emph{core} (domestic-good) inflation, $1+r^{num}=(1+i_{-1})/(1+\pi^H)$.
Second, output $y$ is a physical quantity of the domestic good and is therefore
already expressed in numeraire units---no re-deflation is needed to read it as a
real quantity; CPI-deflated GDP would be $p_{H/P}\,y$, and the two differ only
by the (energy-driven) movement in $p_{H/P}$.

\subsection{The durable state}

Each household carries two binary durable indicators: the technology it
operates this period, $g\in\{B,G\}$, and the technology it held last period,
$g_{-}\in\{B,G\}$. The switching cost $\psi_g$ is incurred exactly in the
transition $g_{-}=B\to g=G$---a household that is green now but was brown last
period. Brown is absorbing; a green durable breaks down to brown at rate
$\delta_g$, so a permanent green share is sustained by a flow of switchers.
These are \emph{two} states, each taking two values; \emph{computationally}
we collapse the pair $(g,g_{-})$ into a single four-valued index, because the
sequence-space household block carries one durable dimension. Of the four
combinations, $(g,g_{-})=(B,G)$---brown now, green last period---never occurs on
the equilibrium path: no household abandons the cheaper technology, and green
is lost only through breakdown, which is already recorded as $g_{-}=B$. It
therefore carries zero mass and the effective state is three-point. This is an
artefact of the collapse, not an economic assumption: the model has two binary
technology states, not one four-way categorical state.
Stages, in forward order, are \texttt{[dep, prod, durables, consav]}:
breakdown, idiosyncratic productivity, the discrete adoption choice (a logit
with taste scale $\sigma_{\varepsilon}$), and the EGM consumption--savings
step. The adoption stage uses the discrete-continuous structure of
\citet{DCEGM2017}.

Three subtleties are worth making explicit, because the calibration and the
sign results depend on them. (i) The choice compares continuation values across
durable states; the fiscal transfer enters cash-on-hand \emph{identically}
across durable states, so it shifts adoption only through the curvature of the
value function, which is why the transfer barely moves the green share.
(ii) The asymmetry ``brown absorbing, green depreciating at $\delta_g$'' means
green ownership requires \emph{active} re-adoption each time the durable breaks
down, so a permanent green share is sustained by a flow of switchers
$D^{sw}=\delta_g D^{G}$ in steady state. (iii) The logit taste scale
$\sigma_\varepsilon$ smooths the choice: at exact indifference the model assigns
a structurally fixed switching probability (not one half), a consequence of the
asymmetric choice sets across states. The magnitude of the adoption response is
governed jointly by $\psi_g$ and $\sigma_\varepsilon$; see
Section~\ref{sec:limits}.

\subsection{Energy prices}

We distinguish three prices, all relative to the CPI:
\begin{align*}
p^{E} &\quad\text{wholesale/retail market price of brown energy, before any subsidy}\\
P^{E}_{B} &= (1-\tau^{E})\,p^{E}+\tau^{E}\bar p^{E}
&&\text{price paid by a \emph{brown} household}\\
P^{E}_{G} &= \varrho\,\bar p^{E}\left(P^{E}_{B}/\bar P^{E}_{B}\right)^{\rho_g}
&&\text{price paid by a \emph{green} adopter}
\end{align*}
with $\varrho$ the steady-state green/brown ratio and $\rho_g$ the
pass-through of the fossil shock to the green price. In the baseline
$\rho_g=0$, so $P^{E}_{G}$ is constant: adopters are fully insulated.

\subsection{How the durable enters the budget}

Households of different durable types face different energy prices, hence
different \emph{consumption bundle prices}. We give each type its own exact
CES cost-of-living index, reusing ARS's own energy-nest parameters
($\alpha_E$, $\eta_E$):
\[
p^{E}_{d}=P^{E}_{B}+\mathbb{I}[d\ \text{green}]\,(P^{E}_{G}-P^{E}_{B}),
\qquad
p_{d}=\Big[1+\alpha_E\big((p^{E}_{d})^{1-\eta_E}-(P^{E}_{B})^{1-\eta_E}\big)\Big]^{\frac{1}{1-\eta_E}},
\]
and the budget constraint is $(p_{d}/p^{num})\,c+a'=\mathrm{coh}(d)$. The
second expression is exact given the CES identity enforced by the equilibrium
block, so brown households have $p_d\equiv 1$ at every date. With
type-specific prices, aggregate energy demand is $\sum_d c_E(d)$, not the CES
demand evaluated at the aggregate index---the two differ by Jensen's
inequality (0.5\% here)---so aggregate $c_H$, $c_F$, $c_E$ are built from
the household block's own heterogeneous-agent outputs.

\subsection{What the CPI measures}
\label{sec:cpi}

With two energy prices in the economy, the energy leg of the CPI could be
anchored on the brown price alone or on the share-weighted index
$[(1-D^{G})(P^{E}_{B})^{1-\eta_E}+D^{G}(P^{E}_{G})^{1-\eta_E}]$. We anchor on
the brown price: it is what statistical agencies do (the energy basket is not
reweighted quarterly for technology adoption), it preserves $p_{d}\equiv 1$
for brown households exactly, and feeding $D^{G}$ back into the CPI would
close a DAG cycle requiring the CPI to be promoted to a model unknown. Because
the unit of account is the domestic good, this choice does not touch the
resource constraint; it affects only the measured deflator and, through it,
the monetary rule. The gap between the share-weighted and brown-anchored CPI
has closed form
\[
\phi^{1-\eta_E}=1+\alpha_E\,D^{G}\Big[(P^{E}_{G})^{1-\eta_E}-(P^{E}_{B})^{1-\eta_E}\Big],
\]
computed ex post outside the model. On the price shock without policy it
reaches 0.14 annualised percentage points, about 1.5\% of the amplitude of
measured CPI inflation (0.23 pp, 2.3\%, on the supply shock; Figure~\ref{fig:deflator}).

\begin{figure}[H]\centering
\includegraphics[width=0.7\textwidth]{output/fig8_deflator_gap.png}
\caption{Measurement gap from anchoring the CPI on $P^{E}_{B}$ alone: true
minus measured annualised inflation, computed ex post.}
\label{fig:deflator}
\end{figure}

\subsection{Balance of payments}
\label{sec:bop}

Both objects below are required, not optional. Once households pay a switching
cost and pay two different energy prices, their summed budget constraints carry
resource flows that must have an aggregate counterpart, or Walras' law fails
and the current account and asset market do not clear. The two terms are those
counterparts.

We implement two internally consistent bookings of these flows, selectable by a
single switch, and report both; the choice has a first-order effect on the
adoption channel's output sign (Section~\ref{sec:signmap}), so it is a result,
not a nuisance.

\paragraph{Import booking (baseline).} The \textbf{switching cost} is booked as
an import, $\text{imports}_{dur}=\psi_g D^{sw}$: booking it as domestic demand
here creates a marginal dividend owned by nobody, since with $\mu_{ss}>1$ and
capitalised profits extra demand for the home good is not competitively
supplied. The \textbf{energy saving} of adopters,
$(P^{E}_{B}-P^{E}_{G})\,C^{G}_{E}$, is booked as a real import saving---cheaper
energy for adopters is treated as a genuine external resource saving, giving the
windfall an external counterpart consistent with asset-market clearing. All
energy, brown and green, is imported at the world brown price and the green
gap is netted off, so green energy remains partly an import.

\paragraph{Domestic booking (green as a domestic industry).} We also model green
explicitly as a domestic sector. Brown energy is imported at the world price;
green energy is supplied by a domestic, competitive, near-zero-marginal-cost
technology---photovoltaic electricity, or the ambient heat a heat pump moves per
unit of electricity---and the green durable is installed by the same domestic
sector at zero profit. Both flows, the green energy bill $P^{E}_{G}C^{G}_{E}$
and the installation $\psi_g D^{sw}$, are domestic value added, rebated to
households lump-sum; only \emph{brown} energy is imported and only brown clears
against the world supply. This resolves the phantom-dividend problem that
motivates the import convention (the rebate is a competitive, zero-profit
sector, so there is no dividend owned by nobody), and the steady state clears to
machine precision. The two bookings are genuinely non-equivalent: the import
booking leaves green partly in the import bill, the domestic booking removes it
entirely and credits the economy with $P^{E}_{G}C^{G}_{E}$ of domestic value
added, a difference of about $0.16\%$ of GDP at the steady state. Because green
production is modelled as costless value added (no green investment, no factor
use), the domestic booking is an \emph{upper bound} on the adoption benefit
(Section~\ref{sec:limits}); the robust object is the sign of the channel, which
differs across the two bookings.

\subsection{Policy instruments}

Both instruments are the extreme (full-compensation) versions, as in
\citet{Bayer2026}.
\begin{align}
\text{Price cap:}\quad & P^{E}_{B,t}=(1-\tau^{E})p^{E}_{t}+\tau^{E}\bar p^{E},
&& \tau^{E}=1 \Rightarrow \text{retail price fixed at pre-crisis level}\\
\text{Slutsky transfer:}\quad & tr_{it}=\iota\,(p^{E}_{t}-\bar p^{E})\,\bar c^{B}_{E,i},
&& \iota=1 \Rightarrow \text{full compensation on \emph{pre-crisis brown} use}
\end{align}
Both are debt financed and repaid through the distortionary labour tax. The
transfer base $\bar c^{B}_{E,i}$ is household $i$'s steady-state \emph{brown}
energy demand, not its total energy demand: the shock raises only the brown
price ($P^{E}_{G}$ is fixed under $\rho_g=0$), so the exact Slutsky compensation
is indexed on brown consumption alone; indexing on total energy would
over-compensate the $5\%$ steady-state green mass, whose price does not rise.
Crucially the transfer is a genuine \emph{lump sum}: the base is the pre-crisis
brown bill and it is paid identically whatever durable the household currently
holds, so it leaves the current brown--green relative price intact. Making it
durable-state-contingent instead---paying only households currently brown, and
so nothing to a household that switches---would tax the switching margin and
spuriously shut the adoption channel down; it is precisely the lump-sum indexing
that lets the transfer preserve adoption (Section~\ref{sec:transfer}) where the
cap eliminates it. Aggregated over the stationary distribution the base equals
the economy's pre-crisis brown energy bill, so the government's outlay matches
what households receive and the current account clears.

The cap's fiscal base is booking-dependent. Only the brown price is capped, so
under the domestic booking, where green is a domestic industry, the subsidy is
indexed on brown consumption alone; under the import booking, where the green
part of the subsidy has an import counterpart, it is indexed on total
energy---the same brown/total distinction as the transfer, reflecting the
$0.16\%$-of-GDP green import content that separates the two bookings
(Section~\ref{sec:bop}). The cap operates directly on $P^{E}_{B}$, exactly the
price entering the brown--green gap, and therefore compresses the adoption
incentive.

\section{Calibration}

\begin{table}[H]\centering
\small
\begin{tabular}{lll}
\toprule
Parameter & Value & Source / target \\
\midrule
$\alpha_E$ energy expenditure share & 0.04 & \citet{ARS2024} \\
$\eta_E$ energy substitution & 0.10 & \citet{ARS2024} \\
$r$, $\mu_{ss}$ & 0.01, 1.03 & \citet{ARS2024} \\
$\beta$ heterogeneity (3 types) & spread 0.06 & fixed; implied MPC $\approx0.30$ \\
\midrule
$\delta_g$ green breakdown rate & 0.05 & assumption (5\,yr life) \\
$\varrho=P^{E}_{G}/P^{E}_{B}$ steady-state gap & 0.80 & heat-pump / EV delivered-cost ratio \\
$\psi_g$ switching cost & 0.538 & targets $D^{G}_{ss}=0.05$ \\
$\sigma_{\varepsilon}$ taste scale & 0.05 & calibrated jointly with $\psi_g$ \\
$\rho_g$ pass-through to green price & 0 & upper bound on the channel \\
\bottomrule
\end{tabular}
\caption{Calibration. The steady state is identical across all experiments.
$\psi_g$ and $\sigma_\varepsilon$ are calibrated jointly to the level target
$D^{G}_{ss}=0.05$; $\sigma_\varepsilon$ is not separately pinned down by an
adoption-elasticity moment (Section~\ref{sec:limits}). The gross markup carries
ARS's home-endowment normalisation $\mu_{ss}\!\leftarrow\!\mu_{ss}(1-\zeta^E\alpha_E)$,
so the value entering firm profits is $\approx1.016$; $1.03$ is the
pre-normalisation primitive.}
\end{table}

\paragraph{Robustness to the green--brown cost ratio $\varrho$.} The one
durable-side parameter with a direct empirical counterpart is the steady-state
green/brown cost ratio $\varrho$. We anchor it at $0.80$ from the relative cost
of delivered energy service under adoption: a heat pump with a coefficient of
performance near three delivers heat from roughly a third of the fuel-equivalent
input, which at 2019--21 electricity-to-gas price ratios places the delivered
green cost around $0.7$--$0.8$ of brown; the per-kilometre cost of an electric
against a comparable internal-combustion vehicle is in the same range.
Table~\ref{tab:rho} re-solves the model across $\varrho\in\{0.6,0.7,0.8\}$,
re-calibrating $\psi_g$ each time to hold $D^{G}_{ss}=0.05$. The \emph{ordering}
is invariant: the cap collapses the peak green-share response to essentially
zero at every $\varrho$, the transfer leaves it at its no-policy value, the cap
cushions cumulative output and the transfer boosts it. Only magnitudes move, and
modestly (the no-policy peak green share ranges over $8.4$--$9.2$ pp). The
limiting case $\varrho=1$---green no cheaper than brown---is degenerate for the
level target: no positive switching cost can sustain a $5\%$ green steady state
when adoption confers no cost advantage, so adoption collapses, which is exactly
the sanity check one wants.

\begin{table}[H]\centering\small
\begin{tabular}{lrrrr}
\toprule
$\varrho$ & $\psi_g$ (re-solved) & no policy peak $D^{G}$ & cap peak $D^{G}$ & transfer peak $D^{G}$\\
\midrule
0.60 & 0.714 & 8.39 & 0.03 & 8.53\\
0.70 & 0.624 & 8.81 & 0.02 & 8.95\\
0.80 & 0.538 & 9.18 & 0.01 & 9.32\\
\bottomrule
\end{tabular}
\caption{Robustness to the green--brown cost ratio $\varrho$ (price shock; peak
green-share response, pp over 24 quarters; $\psi_g$ re-solved each row to
$D^{G}_{ss}=0.05$). Cumulative output moves monotonically but the ordering does
not: no policy $-27.9\!\to\!-30.1$, cap $-21.1\!\to\!-20.9$, transfer
$+14.9\!\to\!+12.7$ as $\varrho$ rises from $0.6$ to $0.8$.}
\label{tab:rho}
\end{table}

\section{Experiments and results}

The \emph{price} shock is a 100 log-point rise in the world energy price with
a 16-quarter half-life, as in \citet{ARS2024}. The \emph{supply} shock is a 10\%
cut in available energy for 6 quarters, after which the world price clears
the market endogenously. Impulse responses are reported as $100\times$ the
level deviation from steady state. Variables whose steady state is one (prices,
$w$, $n$, and the population shares $D^{G}$, $D^{sw}$) read directly as percent
or percentage-point deviations, and $C$ (steady state $\approx0.997$) is within
rounding of a percent; $y$ and $\mathrm{gdp}$ have steady state
$\approx0.985$, so their entries are level deviations $\times100$, within about
1.5\% of the exact percentage. Energy quantities are a share of the energy
steady state ($\approx0.040$).

\input{output/tab_summary_core_import.tex}

\subsection{The crisis moves the transition, in a shock-dependent direction}

With no policy, the green share rises by up to 9.2 percentage points under
the price shock (from 5.0\% to 14.2\% of households) and 2.3 pp under the
supply shock (Figure~\ref{fig:baseline}). The output effect of the adoption
margin, however, is not one-signed. Under the \textbf{supply} shock it is
expansionary ($+0.38$ points of cumulative output with no policy, $+1.67$
under the transfer): adopters escape the energy bill and their real income
falls less. Under the \textbf{price} shock it is contractionary ($-1.60$
points with no policy, $-4.65$ under the transfer): the switching cost is
booked as an import in exactly the periods household income is falling
hardest (Section~\ref{sec:bop}), and at the calibrated $\psi_g$ this resource
cost outweighs the energy-bill saving. The sign of the adoption channel's
contribution to output is therefore not a general property of the model---it
depends on the calibrated switching cost, on which shock is driving
adoption, and (Section~\ref{sec:limits}) on the accounting of $\psi_g$.

\begin{figure}[H]\centering
\includegraphics[width=0.92\textwidth]{output/fig1_price_shock.png}
\caption{Brown energy \emph{price} shock (elastic supply), adoption open
(solid) vs.\ shut (dashed), no policy.}
\label{fig:baseline}
\end{figure}

\begin{figure}[H]\centering
\includegraphics[width=0.92\textwidth]{output/fig1_supply_shock.png}
\caption{Brown energy \emph{supply} shock ($-10\%$ for 6 quarters, fixed
quantity), adoption open (solid) vs.\ shut (dashed), no policy.}
\label{fig:baseline_supply}
\end{figure}

\subsection{A sign condition for the adoption channel}
\label{sec:signmap}

The shock-dependence of the adoption channel's output contribution is not
accidental; it follows from a race between two resource flows, which we can
characterise and then verify. Write the contribution to cumulative output over
horizon $H$ as $\Delta_H\equiv\sum_{t<H}\big(y^{\text{adopt}}_t-y^{\text{no-adopt}}_t\big)$.
Adoption moves two things in the resource constraint (Section~\ref{sec:bop}):

\begin{itemize}\itemsep2pt
\item a \textbf{switching cost} $\psi_g D^{sw}_t$, booked as an import---a
resource \emph{outflow}, hence contractionary---which is \emph{front-loaded}
(paid once, up front, by each switcher) and scales with the adoption
\emph{volume} $D^{sw}$;
\item a \textbf{flow energy saving} $(P^{E}_{B}-P^{E}_{G})\,C^{G}_{E,t}$, booked
as reduced imports---a resource \emph{inflow}, hence expansionary---which
accrues \emph{gradually} over the life of the green durable.
\end{itemize}

\noindent The sign of $\Delta_H$ is the sign of the second flow, cumulated over
the window, minus the first. Two comparative statics follow.

\paragraph{Persistence (price closure).} A more persistent brown-price shock
raises the return to switching and therefore the adoption volume $D^{sw}$
steeply (peak $D^{G}$ rises from $1$ to $22$ pp as the half-life goes from $2$
to $256$ quarters). The front-loaded import cost thus grows faster than the
windowed saving, so $\Delta_H$ is \emph{decreasing} in persistence and crosses
zero at a finite half-life---about $11$ quarters in our calibration
(Figure~\ref{fig:signmap}). Short-lived shocks trigger little adoption whose
quick saving dominates ($\Delta_H>0$); persistent shocks trigger a large
adoption wave whose front-loaded cost dominates within the window
($\Delta_H<0$). This is the opposite of the naive intuition that more
persistence, by raising the per-switcher saving, should make adoption more
expansionary: the per-switcher calculus is dominated by the response of the
switcher \emph{count}.

\paragraph{Closure.} Under the fixed-quantity (supply) closure a third,
expansionary term appears that is absent under the price closure: adoption
lowers fossil demand against a fixed supply, so the market-clearing price
falls, relaxing the scarcity for \emph{all} households. This shifts $\Delta_H$
upward, which is why the supply shock's adoption contribution is positive
($+0.38$) at a persistence where the price shock's is negative.

\paragraph{Booking.} Both statics are conditional on the import booking of the
adoption flows. We solve the model under the alternative domestic booking of
Section~\ref{sec:bop}---green energy and green installation as competitive,
zero-profit domestic value added rebated to households, with only brown energy
imported---which resolves the phantom-dividend problem and clears to machine
precision. Under this booking neither the switching cost nor the green energy
bill is a leakage, and the adoption channel's contribution to cumulative output
($\Delta_{24}$) flips sign:
\begin{table}[H]\centering
\small
\begin{tabular}{lcc}
\toprule
$\Delta_{24}$, adoption output contribution & import booking & domestic booking\\
\midrule
price shock & $-1.60$ & $+35.98$\\
supply shock & $+0.38$ & $+5.96$\\
\bottomrule
\end{tabular}
\caption{Adoption channel contribution to cumulative output ($\Delta_{24}$), import vs.\ domestic booking (Section~\ref{sec:bop}). \emph{None} policy.}
\label{tab:signmap}
\end{table}

The adoption wave itself is essentially unchanged across bookings (peak $D^{G}$
$9.18\to9.46$ pp on the price shock, $2.33\to1.34$ on the supply shock): only
the booking of its resource flows moves, not the margin. The domestic-booking
magnitudes are an \emph{upper bound}---green value added is modelled as costless,
maximally stimulative---but the \emph{sign reversal} is robust and confirms the
characterisation: the sign of the adoption channel is a property of how its
resource cost is booked, not of the adoption margin itself. Finally, the marginal switcher's steady-state indifference,
$\psi_g\approx(\bar P^{E}_{B}-\bar P^{E}_{G})\,\bar C^{G}_{E}/(r+\delta_g)$,
ties the threshold to the cost ratio $\varrho$ and reproduces the direction of
the $\psi_g(\varrho)$ mapping in Table~\ref{tab:rho}. A closed-form threshold
via the full Keynesian-cross algebra of the adoption block is left for the next
draft; the characterisation above is verified numerically in
Figure~\ref{fig:signmap}.

\begin{figure}[H]\centering
\includegraphics[width=0.68\textwidth]{output/fig_signmap.png}
\caption{Adoption's contribution to cumulative output, $\Delta_{24}$, against
the brown-price-shock half-life (price closure, orange), with the supply-shock
value marked (blue). The contribution is expansionary for transitory shocks,
crosses zero near an $11$-quarter half-life, and turns increasingly
contractionary as the shock---and hence the adoption wave whose front-loaded
switching cost it imports---becomes more persistent.}
\label{fig:signmap}
\end{figure}

\subsection{The price cap eliminates the adoption margin}
\label{sec:cap}

Under the cap, the peak green share response falls to essentially zero under
both shocks (from $+9.18$ to $+0.01$ pp on the price shock, from $+2.33$ to
$+0.01$ on the supply shock): the adoption channel is not attenuated, it is
switched off. With the margin shut its contribution to cumulative output
collapses toward zero accordingly. The mechanism is direct: the cap holds
$P^{E}_{B}$ at its pre-crisis level, and $P^{E}_{B}$ is precisely what determines
the return to adoption.

\begin{figure}[H]\centering
\includegraphics[width=0.92\textwidth]{output/fig2_policy_price.png}
\caption{Fiscal response to the \emph{price} shock, adoption open: no policy
(solid), price cap (dashed), Slutsky transfer (dotted).}
\label{fig:policy_price}
\end{figure}

\begin{figure}[H]\centering
\includegraphics[width=0.92\textwidth]{output/fig2_policy_supply.png}
\caption{Fiscal response to the \emph{supply} shock, adoption open: no policy
(solid), price cap (dashed), Slutsky transfer (dotted).}
\label{fig:policy_supply}
\end{figure}

\subsection{The transfer preserves adoption, at a cost to output}
\label{sec:transfer}

The Slutsky transfer leaves the green share response essentially unchanged
($9.46$ vs.\ $9.18$ pp under the price shock, $3.35$ vs.\ $2.33$ under the
supply shock) while delivering more output stabilisation than the cap, at
similar fiscal cost: subsidies destroy the transition margin, transfers do
not. This preservation is a direct consequence of the transfer's lump-sum
indexing (Section~\ref{sec:bop}): because it does not condition on the
household's current durable, switching forfeits nothing and the relative-price
incentive to adopt is left intact. Leaving adoption intact is not free for output under the price shock,
however: the transfer's cumulative-output gain is \emph{smaller} with
adoption open than shut ($-4.65$ points from the margin, Table~\ref{tab:summary}),
for the same resource-cost reason as above.

\subsection{Fixed supply makes the cap destructive}

Under the supply shock the cap yields cumulative output of $-46.9$ against
$-8.2$ with no policy: with inelastic supply, capping the price prevents
substitution away from energy and pushes the economy towards Leontief. This
survives the addition of an adoption margin. It is also precisely the
configuration in which \citet{Langot2026}, assuming elastic supply, find the
cap effective: our elastic-supply column agrees with them ($-20.9$ vs.\
$-30.1$ for cap vs.\ no policy under the price shock).

\begin{figure}[H]\centering
\includegraphics[width=\textwidth]{output/fig3_policy_adoption.png}
\caption{Policy and the adoption margin. Under both shocks the price cap
(dashed) flattens the green share response to zero; the transfer (dotted)
leaves it intact.}
\end{figure}

\subsection{The cap is a dial, not a switch}

Table~\ref{tab:summary} used $\tau^{E}=1$, a complete price freeze. Table~\ref{tab:dose}
sweeps $\tau^{E}\in[0,1]$. Under the \emph{price} shock the adoption loss is
almost exactly linear in $\tau^{E}$ (half a cap costs a bit under half the
transition: 4.53 against 9.18 pp). Under the \emph{supply} shock it is convex:
small caps barely touch adoption and the collapse is concentrated near the
full freeze. A partial shield is therefore much less damaging to the
transition when supply is inelastic---but, as the next result shows, much
more damaging to everything else.

\input{output/tab_dose_core_import.tex}

\subsection{The welfare effect of the cap flips sign with the supply elasticity}

This is the sharpest result in the paper. Under the price shock, tightening
the cap \emph{raises} welfare monotonically, from $-1.30\%$ to $-0.57\%$;
under the supply shock it \emph{lowers} welfare monotonically, from $-0.59\%$
to $-1.28\%$. The sign of the cap's welfare effect therefore hinges on the
supply elasticity. \citet{Langot2026}, assuming elastic supply, find the French
shield effective; \citet{Bayer2026}, assuming fixed union-wide supply, find
subsidies destructive. Our two closures reproduce the \emph{sign} each obtains,
which locates one source of the contrast in a single parameter. We do not claim
to reconcile the two papers fully: Langot's result also runs through a
labour-cost channel and Bayer's through cross-border spillovers in a monetary
union, and our small open economy nests neither. The narrower point we
establish is that, once the supply elasticity is fixed, the direction of the
cap's welfare effect is determined---so part of the empirical disagreement
reduces to measuring that elasticity at the relevant horizon.

\begin{figure}[H]\centering
\includegraphics[width=\textwidth]{output/fig6_dose_response.png}
\caption{Dose-response in $\tau^{E}$ (orange) against the Slutsky transfer
benchmark (dotted). Note the sign reversal of the welfare panel between the
two shocks.}
\end{figure}

\paragraph{Private welfare and the transition are distinct.} The CEV above is a
\emph{private} household-welfare measure: it values consumption and leisure
along the transition and prices no climate or transition externality. Two
results should therefore be kept separate. First, on private welfare alone the
Slutsky transfer dominates the cap under \emph{both} shocks ($-0.25$ vs.\
$-0.57$ under the price shock, $-0.16$ vs.\ $-1.28$ under the supply shock,
Table~\ref{tab:dose}); the case for the transfer over the cap does not rest on
the transition. Second, and separately, the transfer preserves the green
adoption margin that the cap eliminates (Section~4.3)---a positive statement
about the energy transition, not a welfare ranking. We deliberately do not
aggregate the two into a single social-welfare number: that would require a
social value of adoption (a carbon price times avoided brown-energy use) that
we do not model. A positive such value only reinforces the private ranking; it
cannot reverse it here, since the transfer already dominates.

\subsection{Distributional incidence}

\input{output/tab_cev_core_import.tex}

Without policy the impatient lose far more than the patient ($-2.10$ vs.\
$-0.14$, roughly $15\times$). The patient remain the least-hurt group
throughout the cap sweep. One clear distributional result: under the
transfer, the \textbf{patient type's CEV turns positive} ($+0.17\%$) while
impatient and middle types remain negative. Indexed to pre-crisis energy use,
the Slutsky compensation over-compensates low-MPC, low-energy-exposure
households relative to their true loss---a symptom of the fiscal block's
lack of investment margin to crowd out the transfer (Section~\ref{sec:limits}).

\section{Relation to the literature}
\label{sec:lit}

This section is written for the coauthors: it states, paper by paper, what our
setup changes relative to the reference and what that change does to the
results. Compact one-page summaries of each reference are in
Appendix~\ref{sec:summaries}.

\subsection{Relative to \citet{ARS2024}}

\paragraph{Setup.} We inherit the entire ARS aggregate economy---the CES
energy nest, sticky import prices, the mutual fund, the energy suppliers with
intertemporal arbitrage, the real-wage-stabilising wage Phillips curve, and the
constant-real-rate monetary baseline---and change the household in four coupled
ways. (i) ARS's household is homothetic with a fixed energy technology: every
household shares one CES basket and one CPI, and energy demand is
$\alpha_E(P_E/P)^{-\eta_E}c$. We add a discrete brown/green durable state and a
DC-EGM adoption choice, so the energy technology becomes an endogenous,
household-level object. (ii) ARS has a single household energy price; we have
two ($P^E_B$, $P^E_G$) and a type-specific cost-of-living index. (iii) ARS uses
the CPI as unit of account; we use the domestic good, so that the resource
constraint is invariant to how the CPI weights brown and green energy. (iv) We
add two current-account terms absent in ARS---the switching cost as an import,
the adopters' energy saving as an import saving. The energy-supply closure is a
generalisation of ARS: their baseline is our $\varepsilon^E=\infty$
price-taking case, and their coordinated world equilibrium (endogenous world
price) is the two-country analogue of our $\varepsilon^E$-finite case.

\paragraph{Results.} ARS's central mechanism---low $\chi$ makes the energy
price shock contractionary through the real-income channel---is preserved and
operates in the background of every one of our experiments. Where we depart is
that ARS's fiscal ranking is entirely about output, inflation and inequality:
in ARS the subsidy is a domestic ``silver bullet'' that stabilises all three,
and its only drawback is a \emph{cross-country} externality (coordinated
subsidies spike the world price and become self-defeating). Our model adds an
orthogonal cost of the same subsidy that ARS cannot see: capping $P^E_B$
eliminates the green adoption margin (peak $D^G$ from 9.18 to 0.01 pp), so the
subsidy is not neutral for the transition. Two further connections are worth
flagging to the coauthors. First, our supply-shock result---the cap is
destructive under inelastic supply---is the \emph{single-country} counterpart
of ARS's coordination result: in both, capping the price under an
endogenous/scarce energy price prevents the necessary substitution and deepens
the recession. Second, our monetary baseline being ARS's constant-real-rate
rule is what pins the real exchange rate, so a domestic transfer has no
terms-of-trade effect; this is inherited, not assumed anew.

\subsection{Relative to \citet{Bayer2026}}

\paragraph{Setup.} Bayer et al.\ is a two-country HANK monetary union (Germany
one-third, rest of the union two-thirds), with a fixed nominal exchange rate,
built to study \emph{spillovers}. Our economy is a single small open economy
with an exogenous foreign block, so by construction we cannot represent their
central object---the cross-border, beggar-thy-neighbour transmission of a
national subsidy. Three further setup differences matter. (i) Supply: Bayer
fix the union-wide energy \emph{quantity} and let a common price clear; this is
exactly our $\varepsilon^E$-finite closure, so our supply-shock experiments are
the natural comparison. Our shock is smaller (a 10\% cut for 6 quarters against
their 20\% cut for 6 quarters), chosen for first-order reliability. (ii) Supply
side: Bayer has a rich medium-scale block---two assets (liquid bonds, illiquid
capital), an investment margin with capacity utilisation, energy in production
as well as consumption, and GHH preferences. We have a single asset, no
capital, and energy in consumption only in the baseline. (iii) Heterogeneity:
Bayer's headline distributional mechanism is \emph{within-country
energy-share heterogeneity}---poor households have systematically higher energy
shares (a high/low-intensity Markov type correlated with income), so they
suffer larger real-income losses. We assume homothetic energy demand
conditional on the durable type, so we cannot speak to their energy-exposure
distributional results; our heterogeneity is over discount factors and over the
adoption margin instead. Finally, their monetary policy is a standard Taylor
rule, not the constant-real-rate rule; they show this matters, because under a
Taylor rule the subsidy also suppresses \emph{intertemporal} substitution
through household-specific real rates.

\paragraph{Results.} Under the supply closure we reproduce the qualitative
core of Bayer et al.: the Slutsky/Hicks transfer welfare-dominates the price
cap (our private CEV $-0.16$ vs.\ $-1.28$), and capping the price under a fixed
quantity is inefficient because it prevents the substitution the scarcity
requires---our ``cap pushes towards Leontief'' is their ``subsidising under a
supply contraction is inefficient.'' What we add is the adoption margin, which
Bayer et al.\ do not have: the same cap that is inefficient for output also
switches off the transition, and the transfer that is welfare-superior also
preserves it. What we \emph{cannot} deliver is their spillover result: with no
second country, we have nothing to say about whether one country's subsidy
harms its neighbours. This is the reason we are careful (Section~4.6, and the
abstract) not to claim that a single parameter ``reconciles'' Bayer and Langot;
we reproduce the domestic sign, not the cross-border mechanism.

\subsection{Relative to \citet{Langot2026}}

\paragraph{Setup.} Langot et al.\ study France as a single economy inside the
euro area, with an ECB Taylor rule down-weighted for France's size. Two
structural choices distinguish their model from ours. (i) They take the energy
price as \emph{exogenous} (fed from the government's official forecasts), so
France is a price taker facing a perfectly elastic supply at that price---this
is our $\varepsilon^E=\infty$ closure, and their price-shock scenario is the
right comparison. (ii) Their energy substitution elasticity is $\eta_E=0.5$,
five times ours, and they add an \emph{incompressible} (subsistence) level of
energy consumption, which delivers non-homotheticity: poor households have a
larger incompressible share, hence a higher effective energy share and a lower
price elasticity. This is their distributional mechanism, analogous in spirit
to Bayer's energy-share heterogeneity but generated by subsistence rather than
by types. Our model buys energy substitution on a different margin: we keep the
\emph{intensive} elasticity low ($\eta_E=0.1$, as in ARS) and obtain
substitution from the \emph{extensive} adoption margin instead. A third,
methodological difference matters for how the results should be read: Langot et
al.\ conduct an \emph{ex-ante, conditional-forecast} evaluation---they back out
the sequence of structural shocks that makes the HANK reproduce the French
government's published forecasts for output, inflation and debt, then run
counterfactual policies holding those shocks fixed. Ours (like ARS and Bayer)
is a clean MIT-shock impulse-response exercise, not a fit to a specific
historical episode.

\paragraph{Results.} Under the elastic (price) closure we agree with Langot's
headline: the cap cushions output ($\sum y = -20.9$ vs.\ $-30.1$ with no
policy) and, on private welfare, tightening the cap raises the CEV
($-1.30\to-0.57$)---consistent with their finding that the shield supports
growth and curbs inflation. But two qualifications are important for the
coauthors. First, Langot's pro-shield result rests substantially on a
\emph{labour-cost} channel: by containing the retail energy price the shield
prevents the price--wage spiral, holds down the cost of labour, and thereby
supports employment; in their comparison the shield beats both wage indexation
and a targeted transfer precisely on this employment margin. Our wage block and
our lack of an investment/employment margin mean we reproduce the \emph{sign}
of their result but not the full mechanism---hence we do not claim the
Langot--Bayer contrast is ``only'' about the supply elasticity. Second, Langot
et al.\ hold the energy technology fixed, so like ARS and Bayer they do not see
the cap's cost to the transition, which is our contribution.

\section{Limitations}
\label{sec:limits}

\begin{enumerate}\itemsep3pt

\item \textbf{The taste scale $\sigma_{\varepsilon}$ is calibrated, not
estimated.} $\psi_g$ and $\sigma_\varepsilon$ jointly hit $D^G_{ss}=0.05$, but
one level moment cannot separately identify both a switching cost and a taste
scale: a second moment (e.g.\ the semi-elasticity of adoption to a purchase
subsidy) is needed and would discipline the sign found in
Section~4.1 along with the magnitudes throughout. The \emph{ordering} of the
cap-vs-transfer results does not depend on it, since it follows from the
relative price alone.

\item \textbf{The switching cost is booked as an import by accounting
necessity.} Under ARS's $\mu_{ss}>1$ with capitalised profits, booking
$\psi_g D^{sw}$ as domestic demand creates a marginal dividend owned by nobody
and breaks asset-market clearing; we therefore book it externally
(Section~\ref{sec:bop}). This convention is what makes the adoption channel
\emph{contractionary} under the price shock, since the cost lands in the
periods household income is falling hardest (Section~4.1). The \emph{ordering}
of the policy results does not depend on it---it follows from the relative
price alone---but the sign of the adoption channel's output contribution under
a price shock does, and should be read as conditional on this booking.

\item \textbf{Zero pass-through to the green price} ($\rho_g=0$) means
adopters are completely insulated from the fossil shock: an upper bound on
the adoption channel. A partial pass-through calibrated to observed
electricity--gas co-movement would be more defensible.

\item \textbf{No modelled green technology.} The adopters' resource saving
rests on the domestic/imported interpretation of Section~\ref{sec:bop} rather
than on a green sector with its own capacity, cost curve and investment. We
have implemented the minimal version---a competitive, zero-profit installer
sector that books the switching cost as domestic value added rebated to
households---and it clears to machine precision and flips the sign of the
adoption channel (Section~\ref{sec:signmap}); what remains unmodelled is the
green energy \emph{supply} technology (capacity, cost curve, investment) and a
distribution of installer income more realistic than the lump-sum rebate.

\item \textbf{No within-country energy-share heterogeneity.}
\citet{Bayer2026}'s central distributional mechanism is that poor households
have higher energy shares; we have homothetic energy demand conditional on
the durable type, so we cannot speak to their distributional results.

\item \textbf{No investment margin.} Both \citet{Bayer2026} and
\citet{ARS2024} have richer supply sides. Its absence here likely makes the
transfer look more expansionary, and the supply shock more inflationary, than
it would be with capital to smooth the adjustment.

\end{enumerate}

\appendix

\section{Summaries of the three reference papers}
\label{sec:summaries}

These summaries are in our own words, for internal reference. They are
deliberately compact---roughly one page each---and emphasise the features that
matter for our comparison in Section~\ref{sec:lit}.

\subsection{\citet{ARS2024}, ``Managing an Energy Shock''}

\citet{ARS2024} ask how an energy-importing economy should respond to an energy
price shock, and whether household heterogeneity changes the answer. The model
is a small open economy HANK built on the exchange-rate model of Auclert,
Rognlie, Souchier and Straub (2021), extended with an energy good. There are
three goods: a domestically produced home good, an imported foreign good, and
imported energy. Households face idiosyncratic productivity risk, cannot insure
it, and hold a single asset through a mutual fund; preferences are CRRA in a
CES consumption basket that nests energy against a home/foreign bundle, with a
low energy substitution elasticity ($\eta_E=0.1$) and a total substitution
index $\chi=0.3$. Crucially, all households share the \emph{same} basket and
the \emph{same} price index---the energy technology is fixed and homothetic.
Wages are sticky with a real-wage-stabilisation motive; the monetary baseline
is a rule that holds the real interest rate constant. Energy is supplied by
price-taking foreign firms with an intertemporal extraction margin, so in the
single-country equilibrium the world energy price is exogenous.

The analytical core is a decomposition of the output response into an
expenditure-switching channel (positive) and a real-income channel (negative)
scaled by the matrix of intertemporal MPCs. In the complete-markets
representative-agent benchmark only expenditure switching survives, so the
shock is expansionary; with realistic MPCs and low $\chi$ the real-income
channel dominates and the shock is contractionary---their central result, and
the one we inherit. A neutrality result holds at $\chi=1$, where the two
channels cancel. With real wage rigidity a wage--price spiral can appear, but
the real wage still falls.

On policy, ARS study monetary tightening and three fiscal tools---energy
subsidies, targeted transfers (indexed to counterfactual energy use), and
untargeted transfers, all deficit-financed and repaid by a labour tax. For a
country acting alone, subsidies are the most effective domestic stabiliser and
tame inflation (they cap the price that drives the real-wage loss); transfers
support demand but are more inflationary. The decisive twist is
\emph{cross-country}: monetary tightening has a positive externality (it lowers
world energy demand and price), while subsidies have a large negative
externality (coordinated subsidies make world energy demand nearly inelastic,
spike the world price, and become self-defeating). Their policy recommendation
is coordinated tightening plus transfers targeted to the poor, avoiding
subsidies. They also show the impulse responses are not strongly nonlinear at
empirically relevant shock sizes.

For our purposes the key features are: fixed homothetic energy technology, a
single energy price and CPI, the constant-real-rate monetary rule, and the fact
that the only ``cost'' of a subsidy in their world is a cross-country
externality---there is no transition margin.

\subsection{\citet{Bayer2026}, ``Hicks in HANK'' / ``Redistribution Within and
Across Borders''}

\citet{Bayer2026} evaluate the fiscal response to the European energy crisis in
a two-country HANK model of a monetary union, calibrated with Germany as
``Home'' (one-third of the union) and the rest of the area as ``Foreign''
(calibrated to Italy). The model is a medium-scale, two-asset HANK: households
hold a liquid union-wide bond and illiquid capital, face idiosyncratic income
risk, and have GHH preferences; there is an investment margin with capacity
utilisation. Energy enters both production (a CES bundle with the
capital--labour aggregate, elasticity 0.2) and household consumption (a CES
bundle with the physical good, elasticity 0.1). A distinctive and, for them,
central feature is within-country heterogeneity in the \emph{energy share}:
households are of high- or low-energy-intensity type, following a persistent
Markov process correlated with income, so that poorer households spend a larger
share on energy and suffer larger real-income losses. Monetary policy is a
standard union-wide Taylor rule; each national fiscal authority taxes labour
and profits and stabilises its debt.

The defining assumption is that the total union energy supply is
\emph{inelastic}: import capacity is fixed in the short run, so a fixed quantity
is set exogenously and a common price clears the market. The crisis is a 20\%
cut in supply for six quarters, in response to which the energy price rises
roughly fivefold, inflation rises about four points, and output falls around
one percent, with consumption falling most for poor, high-energy households.

They compare two policies, each implemented in Home only so that spillovers can
be measured: a subsidy that caps the retail energy price at its pre-crisis
level, and a Hicks/Slutsky transfer indexed to pre-crisis energy consumption
(to households and firms). The subsidy is the more effective domestic
stabiliser but is fiscally expensive and \emph{beggar-thy-neighbour}: by
holding Home's energy demand up it pushes the common price higher, amplifying
the crisis in Foreign. The transfer stabilises less at Home but generates
essentially no spillover and is cheaper. On welfare (consumption-equivalent
variation), the subsidy \emph{reduces} Home welfare---subsidising energy under a
supply contraction is inefficient because it suppresses both intratemporal and
intertemporal substitution---while the transfer \emph{raises} Home welfare, its
cost being mostly the distortion of the taxes that finance it. Their headline
contrast with ARS (subsidy more stabilising than transfer for output) they
attribute to the Taylor rule and the two-country structure, versus ARS's
constant-real-rate rule.

For our purposes the key features are: fixed union-wide quantity (our inelastic
closure), the two-country spillover mechanism that we cannot represent, the
within-country energy-share heterogeneity that we do not have, and the
substitution-preserving property of the transfer, which we share.

\subsection{\citet{Langot2026}, ``The Macroeconomic and Redistributive Effects
of the French Tariff Shield''}

\citet{Langot2026} evaluate the French \emph{bouclier tarifaire}---the 2022--23
cap on gas and electricity prices---in a HANK model of France as a single
economy within the euro area. The model extends Auclert, Rognlie, Souchier and
Straub with energy in both household consumption and production. Households face
idiosyncratic productivity risk and hold a single asset; preferences are CRRA
with a CES energy nest, and the energy substitution elasticity is $\eta_E=0.5$,
markedly higher than in ARS or Bayer. A distinctive feature is an
\emph{incompressible} (subsistence) level of energy consumption: only energy
above the subsistence level yields utility. This generates non-homotheticity
without non-homothetic preferences---poorer households devote a larger share to
incompressible energy and have a lower price elasticity of energy demand, which
is the source of the distributional results. The ECB Taylor rule is
down-weighted to reflect France's small weight in euro-area inflation. In the
policy simulations the energy price is taken as exogenous (from the
government's forecasts), so France is effectively a price taker facing
perfectly elastic supply.

The methodology is distinctive and worth flagging: rather than clean MIT-shock
impulse responses, the paper performs an \emph{ex-ante, conditional-forecast}
evaluation. Using the sequence-space Jacobian, the authors back out the
sequence of structural shocks (energy price, government spending, transfers,
plus latent demand, markup and fiscal shocks) that makes the model reproduce
the French government's official published forecasts for output, inflation and
the debt ratio through 2027. Policies are then evaluated as counterfactuals
that hold this shock sequence fixed and change one instrument.

The main finding is that the tariff shield supports growth (about 1.9\% average
over 2022--23 against 1.0\% without it) and curbs inflation (5.6\% against
7.2\%), at a fiscal cost near 2\% of GDP per year and a modest rise in the debt
ratio, while containing---though not eliminating---the rise in consumption
inequality even though it is untargeted. The mechanism behind its effectiveness
is a \emph{labour-cost} channel: by holding down the retail energy price the
shield prevents the price--wage spiral, contains the cost of labour, and
thereby supports employment and output. On this macro criterion the shield
beats both a faster wage indexation (more inflationary, worse for output) and a
targeted transfer to finance incompressible energy consumption; the targeted
transfer, however, reduces inequality by more.

For our purposes the key features are: exogenous/elastic energy price (our
price-taking closure), the labour-cost channel behind the pro-shield result
(which our wage block only partly captures), the higher intensive substitution
elasticity, and again a fixed energy technology with no transition margin.

\begin{thebibliography}{9}

\bibitem[Auclert et al.(2024)]{ARS2024}
Auclert, A., H. Monnery, M. Rognlie, and L. Straub (2024).
``Managing an Energy Shock: Fiscal and Monetary Policy.''
In \emph{Heterogeneity in Macroeconomics: Implications for Monetary Policy},
Series on Central Banking, Analysis, and Economic Policies, Vol. 30,
Central Bank of Chile, 39--108.

\bibitem[Bayer et al.(2026)]{Bayer2026}
Bayer, C., A. Kriwoluzky, G. J. M\"uller, and F. Seyrich (2026).
``Redistribution Within and Across Borders: The Fiscal Response to an Energy
Shock.'' \emph{Journal of Political Economy Macroeconomics}, forthcoming.
Earlier version: ``Hicks in HANK,'' CEPR DP 18557.

\bibitem[Iskhakov et al.(2017)]{DCEGM2017}
Iskhakov, F., T. H. J{\o}rgensen, J. Rust, and B. Schjerning (2017).
``The Endogenous Grid Method for Discrete-Continuous Dynamic Choice Models with
(or without) Taste Shocks.'' \emph{Quantitative Economics} 8(2), 317--365.

\bibitem[Langot et al.(2026)]{Langot2026}
Langot, F., S. Malmberg, F. Tripier, and J.-O. Hairault (2026).
``The Macroeconomic and Redistributive Effects of the French Tariff Shield.''
\emph{Journal of Political Economy Macroeconomics}, forthcoming.
Earlier version: CEPREMAP Working Paper 2305.

\end{thebibliography}

\end{document}
